%% pc-modulation — filtrage puis modulation d'amplitude, lus sur les spectres.
%% Trois spectres empilés (amplitude en fonction de la fréquence) reliés par deux flèches :
%%   1. spectre initial : fondamental + harmoniques décroissants, et la porteuse très à droite ;
%%   2. après filtrage : seules les raies de plus grande amplitude sont conservées ;
%%   3. après modulation AM : le spectre est TRANSLATÉ autour de la porteuse (cinq raies).
%% Les fréquences ne sont pas à l'échelle : f_p est en réalité très supérieure à f_1.
\documentclass[tikz,border=4pt]{standalone}
\usepackage{liaisons}
\usetikzlibrary{decorations.pathreplacing}
\begin{document}
\begin{tikzpicture}[x=0.6cm, y=0.70cm,
  axe/.style={-{Stealth[length=5pt]}, line width=0.6pt},
  titre/.style={font=\small\bfseries, inner sep=2pt},
  grad/.style={font=\scriptsize, inner sep=2pt},
  eti/.style={font=\scriptsize, inner sep=2pt, align=center},
  fl/.style={-{Stealth[length=7pt]}, line width=2.2pt, solideA}]

%% ---------- une raie : \raie{fréquence}{amplitude}{couleur} ---------------------
\newcommand\raie[3]{\draw[#3, line width=1.6pt] (#1,0) -- (#1,#2);}
%% ---------- axes d'un spectre : \axes{titre} (repère local) --------------------
\newcommand\axes[1]{%
  \draw[axe] (-0.2,0) -- (11.3,0) node[anchor=west, inner sep=2pt, font=\scriptsize] {$f$ (Hz)};
  \draw[axe] (0,-0.1) -- (0,1.55);
  \node[grad, anchor=south east, inner sep=1pt] at (0.1,1.45) {$u$ (V)};
  \node[titre, anchor=west] at (0.3,1.80) {#1};}

%% =================== 1. spectre du signal à transmettre =========================
\begin{scope}[shift={(0,0)}]
  \fill[solideB!12] (0.55,0) rectangle (6.45,1.15);
  \axes{Spectre initial}
  \foreach \f/\a in {1/1.0, 2/0.72, 3/0.48, 4/0.30, 5/0.18, 6/0.10} { \raie{\f}{\a}{solideB} }
  \raie{9}{1.30}{solideD}
  \node[eti, text=solideB, anchor=south] at (3.6,1.16) {harmoniques};
  \node[eti, text=solideD, anchor=south] at (9,1.32) {porteuse};
  \foreach \f/\lab in {1/{$f_1$}, 2/{$f_2$}, 3/{$f_3$}, 9/{$f_p$}} {
    \draw[line width=0.5pt] (\f,-0.1) -- (\f,0.1);
    \node[grad, anchor=north, inner sep=3pt] at (\f,-0.1) {\lab}; }
\end{scope}

%% ---------------- flèche : filtrage --------------------------------------------
\draw[fl] (5.5,-0.85) -- (5.5,-1.65);
\node[eti, text=solideA, anchor=west, align=left] at (5.9,-1.25)
  {\textbf{filtrage} : on supprime\\les raies de faible amplitude};

%% =================== 2. spectre filtré ==========================================
\begin{scope}[shift={(0,-3.7)}]
  \fill[solideB!12] (0.55,0) rectangle (3.45,1.15);
  \axes{Après filtrage}
  \foreach \f/\a in {1/1.0, 2/0.72, 3/0.48} { \raie{\f}{\a}{solideB} }
  \raie{9}{1.30}{solideD}
  \node[eti, text=solideB, anchor=west, align=left] at (3.7,0.75) {harmoniques conservées};
  \node[eti, text=solideD, anchor=south] at (9,1.32) {porteuse};
  \foreach \f/\lab in {1/{$f_1$}, 2/{$f_2$}, 3/{$f_3$}, 9/{$f_p$}} {
    \draw[line width=0.5pt] (\f,-0.1) -- (\f,0.1);
    \node[grad, anchor=north, inner sep=3pt] at (\f,-0.1) {\lab}; }
\end{scope}

%% ---------------- flèche : modulation ------------------------------------------
\draw[fl] (5.5,-4.55) -- (5.5,-5.55);
\node[eti, text=solideA, anchor=west, align=left] at (5.9,-4.95)
  {\textbf{modulation} : le spectre\\est transporté autour de $f_p$};

%% =================== 3. spectre modulé (AM) =====================================
\begin{scope}[shift={(0,-7.85)}]
  \axes{Après modulation}
  \raie{9}{1.30}{solideD}
  \foreach \f/\a in {7.2/0.36, 8.1/0.52, 9.9/0.52, 10.8/0.36} { \raie{\f}{\a}{solideB} }
  \node[eti, text=black!60, anchor=west, align=left] at (0.3,0.70)
    {plus aucune raie\\au voisinage de 0};
  \foreach \f/\lab in {7.2/{$f_p - f_A$}, 8.1/{$f_p - f_B$}, 9/{$f_p$},
                       9.9/{$f_p + f_B$}, 10.8/{$f_p + f_A$}} {
    \draw[line width=0.5pt] (\f,-0.1) -- (\f,0.1);
    \node[font=\scriptsize, rotate=45, anchor=north east, inner sep=2pt] at (\f,-0.1) {\lab}; }
  %% accolade : bande occupée par le signal émis
  \draw[solideF, line width=0.7pt, decorate,
        decoration={brace, amplitude=4pt, raise=2pt}] (7.2,1.35) -- (10.8,1.35);
  \node[eti, text=solideF, anchor=south] at (9,1.58) {bande de fréquences émise};
\end{scope}

%% =================== encadré de synthèse ========================================
\node[draw, line width=0.6pt, rounded corners=2pt, fill=black!3, inner sep=5pt,
      font=\scriptsize, align=center] at (5.55,-10.70)
  {\textbf{Transposition de fréquence} : le spectre est translaté\\autour de la porteuse, sans perte d'information.};
\end{tikzpicture}
\end{document}
